\id{ҒТАМР 65.09.05}{}

\begin{header}
\swa{А.Ә.~Тоқтабек, Л.А.~Мамаева, Hülya~Gül, А.Е.~Отуншиева, М.Д.~Алибекова}{ЙОГУРТ, БАЛДЫРКӨК ҰНТАҚТАРЫМЕН НЫҒАЙТЫЛҒАН МАКАРОН ӘЗІРЛЕУ БОЙЫНША ЗЕРТТЕУЛЕР}

\tsp{1}А.Ә.~Тоқтабек\envelope,
\tsp{1}Л.А.~Мамаева,
\tsp{2}Hülya~Gül,
\tsp{3}А.Е.~Отуншиева,
\tsp{4}М.Д.~Алибекова
\end{header}

\begin{affil}
\tsp{1}Қазақ ұлттық аграрлық зерттеу университеті, Алматы, Қазақстан,

\tsp{2}Süleyman Demirel University, Isparta, Türkiye,

\tsp{3}Әл-Фараби атындағы Қазақ ұлттық университеті, Алматы, Қазақстан, М.Әуезов атындағы Оңтүстік Қазақстан зерттеу университеті, Шымкент, Қазақстан,

\corrauthor{Корреспондент-автор:~aijan700@mail.ru}
\end{affil}

Макарон өнімдерінің ең көп тұтынатын азық-түлік өнімі екені айтылған.
Қазіргі таңда азық-түлік өнеркәсібінде маңызды бағыттардың бірі болып
саналады. Мұндай өнімдердің дамуы халықтың денсаулығына зиян келтірмеу
үшін дәмді және пайдалы тағамдарды ұсыну мақсат қояды. Макарон өнімдері
пайдалы өнімдерді алу үшін жаңа макарон өнімдерін өндірудегі кейбір
негізгі қиындықтарды қарастырады. Тағамдық макарон өнімдерінің жолдарын
әзірлеу. Макарон өнімдерін өндірудің негізгі шикізаты бидай ұны мен су
болып табылады, ал қосымша шикізатқа әр-түрлі байытқыш қоспалар
(жұмыртқа, йогурт, балдыркөк ұнтақтары) жатады. Тағамдық және
денсаулыққа пайдалы қасиеттері бар макарон өнімдерін жасау технологиясы
қарастырылады. Зерттеу барысында дәстүрлі бидай ұнына йогурт, балдыркөк
ұнтақтары қосу арқылы, функционалдық құрамы байытылған макарон өнімінің
рецептурасы ұсынылды. Йогурт ұнтағы өнімнің ақуыздық құрамын арттырып,
пайдалы микроэлементтермен толықтырды, сондай-ақ өнімнің дәмі мен иісіне
жағымды әсер етті. Алынған макарон үлгісі ГОСТ 14849-89 стандартына сай
пішінін сақтап, қайнатқаннан кейін ыдырамаған және жабыспаған. Сенсорлық
талдау нәтижелері бойынша дайын өнімнің органолептикалық көрсеткіштері
жоғары бағаланды. Зерттеу нәтижесінде денсаулыққа пайдалы әрі тағамдық
құндылығы жоғары макарон өнімін дайындаудың тиімді жолы анықталды.

{\bfseries Түйін сөздер:} макарон өнімдері, қатты бидай, нығайтылған
макарон, тағамдық қасиеттер, йогурт ұнтағы, рецептура, балдыркөк ұнтағы.

\begin{header}
RESEARCH ON THE DEVELOPMENT OF PASTA ENRICHED WITH YOGURT AND CELERY POWDER

\tsp{1}A.A.~Toktabek\envelope,
\tsp{1}L.A.~Mamayeva,
\tsp{2}Hülya~Gül,
\tsp{3}A.E.~Otunshiyeva,
\tsp{4}M.D.~Alibekova
\end{header}

\begin{affil}
\tsp{1}Kazakh National Agrarian Research University, Almaty, Kazakhstan,

\tsp{2}Suleyman Demirel University, Isparta, Türkiye,

\tsp{3}Al-Farabi Kazakh National University, Almaty, Kazakhstan,

\tsp{4}M. Auezov South Kazakhstan Research University, Shymkent, Kazakhstan,

e-mail:~aijan700@mail.ru
\end{affil}

It is said that pasta is the most consumed food product. Nowadays, it is
considered one of the most important directions in the food industry.
The development of such products aims to provide tasty and healthy
dishes without harming the health of the population. Pasta products
consider some of the main difficulties in the production of new pasta
products in order to obtain healthy products. Development of ways of
dietary pasta products. The main raw materials for the production of
pasta products are wheat flour and water, and additional raw materials
include various enriching additives (eggs, yogurt, celery powder). The
technology of making pasta products with nutritional and health benefits
is considered. During the study, a recipe for a pasta product enriched
with functional composition was proposed by adding yogurt and celery
powder to traditional wheat flour. Yogurt powder increased the protein
content of the product, supplemented it with useful microelements, and
also had a positive effect on the taste and smell of the product. The
obtained pasta sample retained its shape in accordance with the GOST
14849-89 standard, did not disintegrate and did not stick after boiling.
According to the results of sensory analysis, the organoleptic
indicators of the finished product were highly appreciated. The research
revealed an effective way to prepare a healthy and nutritious pasta
product.

{\bfseries Keywords\emph{:}} pasta, durum wheat, fortified pasta,
nutritional information, yogurt powder, recipe, celery powder

\begin{header}
ИССЛЕДОВАНИЯ ПО РАЗРАБОТКЕ МАКАРОН, ОБОГАЩЕННЫХ ЙОГУРТОМ И СЕЛЬДЕРЕЕМ ПОРОШКОМ

\tsp{1}А.Ә.~Тоқтабек\envelope,
\tsp{1}Л.А.~Мамаева,
\tsp{2}Hülya~Gül,
\tsp{3}А.Е.~Отуншиева,
\tsp{4}М.Д.~Алибекова
\end{header}

\begin{affil}
\tsp{1}Казахский национальный аграрный, исследовательский университет, Алматы, Казахстан,

\tsp{2}Университет Сулеймана Демиреля, Испарта, Турция,

\tsp{3}Казахский национальный университет им., Аль-Фараби, Алматы, Казахстан, Южно-Казахстанский исследовательский университет им.М. Ауэзова, Шымкент, Казахстан,

e-mail:~aijan700@mail.ru
\end{affil}

Макаронные изделия являются наиболее потребляемым продуктом питания. В
настоящее время они считаются одним из важнейших направлений в пищевой
промышленности. Разработка таких продуктов направлена
на обеспечение вкусных и полезных блюд без вреда
для здоровья населения. Макаронные изделия рассматривают некоторые из
основных трудностей в производстве новых макаронных изделий с целью
получения полезных продуктов. Разработка способов диетических макаронных
изделий. Основным сырьем для производства макаронных изделий являются
пшеничная мука и вода, а дополнительным сырьем являются различные
обогащающие добавки (яйца, йогурт, порошок сельдерея). Рассматривается
технология изготовления макаронных изделий с питательной и
оздоровительной пользой. В ходе исследования была предложена рецептура
макаронного изделия, обогащенного функциональным составом путем
добавления йогурта и порошка сельдерея к традиционной пшеничной муке.
Йогуртовый порошок увеличил содержание белка в продукте, дополнил его
полезными микроэлементами, а также оказал положительное влияние на вкус
и запах продукта. Полученный образец макаронных изделий сохранил форму в
соответствии с ГОСТ 14849-89, не рассыпался и не слипся после варки. По
результатам сенсорного анализа органолептические показатели готового
продукта получили высокую оценку. В ходе исследований выявлен
эффективный способ приготовления полезного и питательного макаронного
изделия.

{\bfseries Ключевые слова:} макаронные изделия, твердые сорта пшеницы,
обогащенные макаронные изделия, питательные свойства, йогуртовый
порошок, рецепт, порошок сельдерея.

\begin{multicols}{2}
{\bfseries Кіріспе.} Қазіргі таңда, әлем бойынша орын алып отырған
бәсекелестік және экономиканың қауіпсіздігін қамтамасыз ету жағдайында
елді жоғары сапалы азық-түлік өнімдерімен қамтамасыз ету мәселесі өте
өзекті болып отыр. Сондықтан тамақ өнеркәсібі кәсіпорындарының
инновациялық дамуы мәселесіне ерекше назар аудару керек, өйткені
инновациялар экономиканың барлық салаларында бәсекелестік тұрғысында
жетістікке жетуінің маңызды болып табылады. Шетелдік бәсекелестік
отандық тамақ өнеркәсібі кәсіпорындарына үнемі өндіріске инновациялар
енгізу арқылы өнімдерді жетілдіру қажеттілігін тудырады. Жаңа өнімдерді
енгізу кәсіпорындарға өзінің пайдасын сақтап қалуға, ал инновациялық
дамуға құйған инвестициялар кәсіпорын шығындарын азайтуға мүмкіндік
береді {[}1{]}. Сонымен қатар, Қазақстан Республикасында макарон
өнімдерін тұтыну Италиядан кейінгі әлемдегі ең жоғары көрсеткіштердің
бірі болып қала береді. Сондықтан, Қазақстан Республикасында қалыптасқан
жағдайға байланысты қазіргі уақытта кәсіпорындар жоғары өнімділікті
сақтап және ресурстарды айтарлықтай үнемдей отырып, өндірілетін өнімнің
сапасына үлкен көңіл бөліп, макарон өнімдерін өндірудің жаңа
технологияларын жасауға көңіл аудару {[}2{]}.

Макарон өнімдерінің тағамдық және әлеуетті денсаулықты жақсартатын
қасиеттерін жақсарту үшін әр-түрлі әдістерді қолдануға болады. Ең
қарапайым әдіс макарон өнімдерін көп жылдар бойы коммерциялық түрде қол
жетімді тұтас дәнді немесе тұтас астық өнімі ретінде тұтыну болуы
мүмкін. Дегенмен, тұтынушылар тазартылған дәндерден немесе ұннан
жасалған макарон өнімдерін жақсы көреді, өйткені олардың тұтас
дәндер/тұтас дәндермен салыстырғанда денсаулыққа пайдасы аз болғанымен,
дәмі, сыртқы түрі және құрылымы жақсы. Макарон өнімдерінің тағамдық
қасиеттерін жақсартудың басқа тәсілдеріне олардың оқшауланған функциясы
туралы білімге немесе басқа тағамдардағы зерттеулерге негізделген арнайы
функцияларды қамтамасыз ету үшін арнайы ингредиенттерді немесе
ингредиенттер комбинацияларын қосу кіреді {[}2{]}. Ресурстарды үнемдеу,
бақылау және оңтайландыру мәселелері алға қойған мақсаттарының бірі,
өйткені олар макарон өнімдерін шығаратын кәсіпорындарда қатты бидай ұнын
шамадан тыс тұтынумен байланысты. Сонымен қатар, макарон өнімдерін
өндірудің технологиялық нормативтік құқықтық базасын жетілдіру және шешу
қажет. Бұл мәселелерді шешу технология және стандарттау салаларындағы
зерттеулерді қажет етеді. Тамақ өнеркәсібі, оның ішінде макарон өнімдері
екі түрлі негізгі құбылыстарға тәуелді болып келеді, өйткені бастапқы
шикізаттың сапасына және оны өңдеудегі технологиялық әрекеттердің
ерекшеліктеріне көңіл бөледі. Макарон өнімдерін өндіруге арналған басты
шикізаттардың біріншісі ұн болып табылады,ал екіншісі су болып саналады.
Ал қосымша шикізаттарға әр-түрлі құнарландырғыш және дәмдеуіш қосымша
заттарды жатқызуға болады {[}3{]}.

Соңғы жиырма жылда тамақтанудың адам денсаулығына әсері туралы хабардар
болу тұтынушылардың тағамдық қасиеттері жақсартылған тағамға деген
қызығушылығын арттырды. Шикізат деп кез-келген өнімді өндіру үшін
бастапқы материал қызметін атқаратын, яғни әрі қарай өңдеуге жататын
адам еңбегінің объектілері түсіндіріледі. Демек, азық-түлік секторы
тағамдық құндылығынан асып түсетін пайдалы қасиеттері бар функционалды
тағамдарды зерттеуге және дамытуға баса назар аударды. Шынында да,
макарон өнімдері өзінің әмбебаптығы, арзан құны, тағамдық құндылығы және
жағымды дәмі арқасында бүкіл әлемде кеңінен тұтынылады, осылайша адам
денсаулығына пайдалы әсер ететін функционалды молекулаларды жеткізу үшін
қолайлы тасымалдаушы болып табылады. Функционалды ингредиенттерді немесе
балама шикізатты пайдалану макарон өнімдерінің пісіру сапасы мен
сенсорлық атрибуттарын жиі өзгертеді, нәтижесінде пісіру жоғалуы артады,
қаттылығы төмендейді, түс өзгереді және макарон дәмі өзгереді. Сондықтан
макарон өнімдерін инновациялау стратегиясы функционалдық ингредиенттерді
анықтаудан басқа тұтынушыларға тартымды болу үшін жоғары пісіру және
сенсорлық сапаны сақтай отырып, пайдалы қасиеттерін сақтай алатын жаңа
макарон өнімдерін алу үшін қолайлы технологиялық процесті қажет етті.
Макарон өнімдерінің функционализациясы негізінен қатты бидай
жармаларында шектеулі мөлшерде болатын диеталық талшықтармен,
антиоксиданттармен, ақуыздармен, витаминдермен және минералдармен байыту
арқылы қол жеткізілді\emph{.} Денсаулығын ойлайтын тұтынушылардың
салауатты тағамға деген өсіп келе жатқан сұранысы азық-түлік
өндірушілерінің қызығушылығын тудырды. Соңғы онжылдықта өндірушілердің
тағамдық құндылығын жақсарту немесе жарманы ішінара немесе толық
ауыстыру ретінде болжамды денсаулық пайдасын қамтамасыз ету үшін
әр-түрлі ингредиенттермен толықтыру үрдісі болды. Ұнға оңай өңделетін
көне дәндер де белгілі бір артықшылықтар бере алады, бірақ олардың
функционалдық құндылығының дәлелі онша анық емес.

Денсаулыққа қосымша пайдасы бар макарон өнімдерін тұтынуды арттыру үшін
зерттеушілер, өнеркәсіп және тиісті агенттіктер сенсорлық сапаны
жақсарту, өңдеу мәселелері (пісіру уақыты), қолжетімділік және аралас
медиа хабарламалары сияқты оларды қабылдаудағы кейбір кедергілерді
еңсеруі керек.

Макарон - танымал тағам және жоғары қоректік немесе функционалды
қосылыстарды қосудың жақсы тәсілі екендігі көрсетілген. Ауыстырылатын
ингредиенттері бар макарон өнімдерінің жақсы технологиялық сапасы мен
тұтынушыға қолжетімді болуын қамтамасыз ету керек. Сондықтан мұндай
өнімдерді өндіруші тиімді болуы керек және өндіруші мұндай өнімдерді
дайындамас бұрын, сондай-ақ әлеуетті нарыққа дейін қажетті ингредиенттің
дайын жеткізілімін қамтамасыз етуі керек. Белсенді қосылыстар мен ақуыз
матрицасы арасындағы өзара әрекеттесу кебек, инулин, еритін талшық және
төзімді крахмал сияқты кейбір ингредиенттер үшін жақсы түсініледі.
Макарон өнімдерінің сапасына әсер ететін жекелеген қосылыстар арасындағы
синергизм немесе өзара әрекеттесу туралы бірнеше зерттеулер ғана
қарастырылған. Дегенмен, көптеген елдердің заңнамасы төмен GI,
холестеринді төмендететін, жүрекке пайдалы және т.б. сияқты тағамға
қатысты денсаулыққа шағым жасамас бұрын дәлелдеуді талап етеді. Осы
себепті денсаулықтың нақты пайдасын растау үшін адам клиникалық
сынақтары арқылы ең перспективалы функционалды макарон өнімдерін бағалау
үшін көбірек зерттеулер қажет. Қол жетімді баға бойынша жақсы дәмі,
құрылымы және сыртқы түрі бар денсаулық туралы мәлімдемелер
тұтынушылардың мұндай тағамдарға сұранысын арттыруға көмектеседі.

{\bfseries Материалдар мен әдістер.} Ғылыми негізделген рецептура бойынша
дайындалған макарон өнімдерінің йогурт ұнтағымен нығайтылған макарон
әзірлеу бойынша зерттеулер үшін сынамалар алынды. Макарон өнімдерін
өндіруге арналған рецептура әзірленді және қамырдың технологиялық
қасиеттері зерттелді. Ұн қоспасының тағамдық құндылығын, реологиялық
қасиеттерін, илеу ұзақтығын және макарон қамырын престеу жылдамдығын
анықтау әдістемесі бойынша химиялық құрамын анықтау (ақуыздың массалық
үлесі, крахмал және талшық мөлшері, күл мөлшері) инфрақызыл
анализатордың көмегімен жүзеге асырылды (1-сурет) {[}3,4{]}\emph{.}

Макарон өнімдеріне йогурттан бес сынама (K1,A1,A2,A3 және A4) құрамында
тек бидай ұны (100\%) бар бақылау ретінде K1 үлгісі қолданылды, ал A1,
A2, A3 және A4 үлгілері әр-түрлі концентрациядағы бидай ұны мен йогурт
ұнтағы арқылы дайындалды, тура сол тәрізді балдыркөк ұнтақтарын қосу
арқылы әр-түрлі концентрациядағы бидай ұны мен балдыркөк ұнтағын қосу
арқылы Б1, Б2, Б3, Б4 және құрамында тек бидай ұны (100\%) бар бақылау
ретінде K1 үлгісі тәрізді 5 сынама дайындалып олардың сапасына
эксперименттер жүргілді. Қатты бидайдан жасалған йогурт және балдыркөк
ұнтақтарын қосу арқылы макарон өнімдерінің тағамдық сапасын жақсарту
үшін орасан зор әлеуетке ие. Барлық сынамалар құрамы 11г жұмыртқа, 35мл
су негізінде жасалынды.

\fig[0.8\columnwidth]{p/image23}[1-сурет. Инфраскан "Кан"]
\vspace{-0.5cm}
Макарон өнімдеріне йогурт ұнтағын және балдыркөк ұнтағын қосу арқылы
сапасын айтарлықтай жақсарта алады. Йогурт ұнтағы-бұл құрғақ күйінде
болатын, дайын болған кезде сумен немесе басқа сұйықтықтармен
араластырылғанда йогурт түрінде дайындалатын өнім. Ол көбінесе өндіру
процесін жеңілдету, сақтау мерзімін ұзартып, тасымалдау мен сақтау
шарттарын жақсарту үшін қолданылады. Балдыркөк ұнтағы- бұл құрғақ
күйінде дайындалған және ұнтақталған балдыркөк өсімдігінің өнімділігі.
Ол көбінесе тағамдарға дәм қосу, дәмдік қасиеттерін арттыру, сонымен
қатар дәрумендер мен минералдарға бай болу мақсатында қолданылады
(2-сурет).

Алынған макарон өнімдерінің үлгілері нормативтік құжаттарда белгіленген
қауіпсіздік талаптарына сәйкес келеді {[}5-9{]} және көрсеткіштер рұқсат
етілген шектік мөлшерден аспайды, яғни бұл өнімдер қауіпсіз деп
бағаланады.

Жүргізілген зерттеулер бойынша макаронның беріктігін (durability) немесе
қаттылығын (hardness) анықтай отырып, осы көрсеткіштер негізінде тағам
өнімінің сапасын, тұтынушы қабылдауын және өндірістік стандарттарға
сәйкестігін бағалауға мүмкіндік аламыз.

Макарон өнімдеріне йогурт ұнтағын және балдыркөк ұнтағын қосу арқылы
дайындалған сынамалар құрамы мен концентрациясы 1 және 2 кестедегідей
мәндерге ие болды.
\end{multicols}

\begin{figs}[2 - cурет. Йогурт және балдыркөк ұнтақтарынның жалпы көрінісі]
    \fig[0.4\textwidth][\textwidth]{p/image24}[Балдыркөк ұнтағы]
    \fig[0.4\textwidth][\textwidth]{p/image25}[Йогурт ұнтағы]
\end{figs}

\tcap{1-кесте. Макарон өнімдеріне йогурт ұнтағын қосу арқылы
дайындалған сынамалар құрамы мен концентрациясы}
\begin{longtblr}[
  label = none,
  entry = none,
]{
  width = \linewidth,
  colspec = {Q[20]Q[100]Q[200]Q[234]Q[44]Q[81]Q[119]Q[95]},
  cells = {c},
  cells = {font = \small},
  cell{1}{2} = {font=\bfseries\small},
  cell{1}{3} = {font=\bfseries\small},
  cell{1}{4} = {font=\bfseries\small},
  cell{1}{5} = {font=\bfseries\small},
  cell{1}{6} = {font=\bfseries\small},
  cell{1}{7} = {font=\bfseries\small},
  cell{1}{8} = {font=\bfseries\small},
  hlines,
  vlines,
}
№ & Сынама № & Құрамын-дағы бидай ұны, \% & Құрамындағы йогурт ұнтағы, \% & Ені & Биіктігі & Пісіру уақыты & Салмағы \\
1 & А1       & 90                         & 10                            & 5,8 & 2,3      & 22 мин 28сeк  & 25,48г  \\
2 & А2       & 85                         & 15                            & 6,1 & 2,0      & 14 мин 18сeк  & 25,13г  \\
3 & А3       & 80                         & 20                            & 6,1 & 2,1      & 11 мин 56 сек & 25,16г  \\
4 & А4       & 70                         & 30                            & -   & -        & -             & -       \\
5 & К1       & 100                        & -                             & 5,8 & 2,3      & 23 мин 43 сек & 25,20г  
\end{longtblr}

\begin{multicols}{2}
Мұнда А4 сынама бойынша қамыр макарон жасауға болмайды. Өйткені оның
икемдігілі және жабысқақтығы өте жоғары.

Зерттеу Stable Micro Systems (Англия) фирмасының текстуралық
анализаторын қолданумен жүзеге асырылды. Stable Micro Systems (Англия)
фирмасының текстуралық анализаторлары қатты заттардың, тұтқыр
сұйықтықтардың, ұнтақтардың және түйіршікті материалдардың реологиялық
қасиеттерін зерттеу үшін іргелі, эмпирикалық және имитациялық сынақтарды
жүргізуге мүмкіндік береді. Сынақтар қысу немесе созу арқылы сынақ
үлгісіне бір реттік немесе циклдік әсер ету арқылы жүзеге асырылады.
Сынақ кезінде деформацияға әсер ету керек күш уақыттың әрбір сәтінде
сынақ аяқталуының белгіленген сәтіне дейін өлшенеді. Нәтижелерді алу
бойынша макарон өнімдеріне сынама жүргізу сәті (3-сурет)
\end{multicols}

\tcap{2-кесте. Макарон өнімдеріне балдыркөк ұнтағын қосу арқылы
дайындалған сынамалар құрамы мен концентрациясы}
\begin{longtblr}[
  label = none,
  entry = none,
]{
  width = \linewidth,
  colspec = {Q[20]Q[100]Q[200]Q[234]Q[44]Q[81]Q[119]Q[95]},
  cells = {c},
  cells = {font = \small},
  cell{1}{2} = {font=\bfseries\small},
  cell{1}{3} = {font=\bfseries\small},
  cell{1}{4} = {font=\bfseries\small},
  cell{1}{5} = {font=\bfseries\small},
  cell{1}{6} = {font=\bfseries\small},
  cell{1}{7} = {font=\bfseries\small},
  cell{1}{8} = {font=\bfseries\small},
  hlines,
  vlines,
}
№ & Сынама № & Құрамын-дағы бидай ұны, \% & Құрамындағы балдыркөк ұнтағы, \% & Ені & Биіктігі & Пісіру уақыты & Салмағы \\
1 & Б1       & 98,5                       & 1,5                              & 5,8 & 2,1      & 18мин 06 сек  & 25,01г  \\
2 & Б2       & 97                         & 3                                & 5,8 & 2,2      & 18 мин 32 сек & 25,21г  \\
3 & Б3       & 95,5                       & 4,5                              & 6,0 & 2,1      & 15 мин 48 сек & 25,10г  \\
4 & Б4       & 94                         & 6                                & 6,0 & 2,0      & 13 мин 25 сек & 25,15г  \\
5 & К1       & 100                        & -                                & 5,8 & 2,3      & 30 мин 43 сек & 25,20г  
\end{longtblr}

\begin{multicols}{2}
\fig[0.8\columnwidth]{p/image26}[3 - cурет. Макарон өнімдеріне сынама жүргізу сәті]
\vspace{-0.5cm}
{\bfseries Нәтижелер және талқылау.} Алынған тәуелділіктер үлгілердің
қаттылығын, серпімділігін, беріктігін, тұтқырлығын, аққыштығын,
консистенциясын және басқа реологиялық параметрлерін бағалауға мүмкіндік
береді.

Бұл құрылғының ерекшелігі - олардың әмбебаптығы (жаңа әзірлемелер үшін
де, халықаралық стандарттарға сәйкес тестілеу үшін де пайдалану
мүмкіндігі) және ең дәл нәтижелерді алуға және үлгілердегі ең аз
айырмашылықтарды алуға мүмкіндік береді. Анализаторда үлгілердің
физикалық және химиялық қасиеттерін бағалау кезінде зертханаларда
қолданудың кең ауқымы үшін қосымшалардың үлкен таңдауы бар.

Макарон өнімдерінің үлгілері (А1, А2, А3 және А4) 100:00 қатынасында
қатты бидай ұны мен йогурт ұнтағының әр-түрлі концентрациясын пайдалана
отырып дайындалды; 90:10; тиісінше 85:15, 80:20, 70:30. Әрбір жағдайда
дайындалған макаронның мөлшері тиісті құрамның 500г болды. Қалыптың
алдыңғы жағы дайын өнімнің ұзындығы әрбір үлгі үшін 4см болатындай етіп
реттелді. Соңғы өнімді кептіру макарон үлгісі пеште 55ºC температурада 1
сағат ұсталды. Кептірілген өнім полиэтилен пакеттерге оралған.
Кептірудің негізгі мақсаты үлгідегі ылғалдылықты шамамен 12-13\% дейін
төмендету болды. Әр-түрлі үлгідегі өнімдер тығыздығы жоғары полиэтилен
пакеттерге оралған. Алынған кептірілген өнімдер пісіру уақыты, суды
талдау, ылғалдылық, құрылым және сенсорлық талдау сияқты қосымша
зерттеулер үшін пайдаланылды {[}10{]}.

Макарон өнімдерінің сапасын зерттеу кезінде тұтынушылық қасиеттерінің
жиынтығы, тағамдық құндылығы және өнімдердің мақсаттарға жарамдылығы
анықталды. Өнімдердің органолептикалық қасиеттері:дәмі, түсі, иісі,
сыртқы түрі. Ылғалдылығын, қышқылдығын, дәмі мен иісін, пісіргеннен
кейінгі макаронның жай-күйін анықтау және макарон өнімдерін қабылдау
және әзірленген макарон өнімдерінің сапа көрсеткіштері бойынша талдаудан
өтті.

Ылғалдың массалық үлесі (W) формула арқылы есептеледі:

\begin{equation}
W = \frac{\left( m_{1} - m_{2} \right)}{m}*100
\end{equation}

Мұндағы: m\tsb{1}- кептіру алдында талдауға арналған сынама
бөтелкесінің салмағы, г;

\(m_{2}\) - кептіргеннен кейін талдауға арналған сынама бөтелкесінің
салмағы, г;

\(m\)- талдауға арналған үлгінің массасы, г

\begin{equation}
W = \frac{{(1,4669}^{\ } - {0,7715}^{\ })\ }{0,7976} \times 100 = 87,186
\end{equation}

\begin{equation}
W = \frac{{(1,4189}^{\ } - 0,7694)\ }{0,7542} \times 100  88,117
\end{equation}

(2) формула бойынша келтірілген 98,5 г ұн, 1,5г балдыркөк ұнтағы және
(3) формула бойынша 90 г ұн, 10 г йогурт ұнтағын қолданылды.

Ылғалдылық тұрақты массаға дейін кептіру әдісі арқылы анықталып, есептеу
осы формула бойынша жүргізілді {[}11{]}.

Макарон өнімдерінің пісіру қасиеттері пісірілгенге дейін пісіру
ұзақтығына, пісіру кезінде сіңірілген судың мөлшеріне, қайнаған суға
өткен құрғақ заттардың мөлшеріне, пісірілген бұйымдардың пішінінің
сақталуына қарай бағаланды. Пісірудің ұзақтығы дайын болғанға дейін
қайнаған суға салынғаннан бастап ұнды, пісірілмейтін қабат жоғалып
кеткенге дейінгі уақыт аралағында болып саналады.Ыдыстан өнімнің бір
бөлігін алып, екі шыны әйнектің арасына қою керек.Макарон өнімдерінің
пісіру аяқталғаннан кейін макарон ыдыстың үстінде ұсталған сүзгіге
ауыстырылды. Артық су ағып болғаннан кейін, сыртқы тексеру арқылы
бұйымдардың пішінін сақталуы белгіленді.

Үлгілердің сапасы AACC 2000 стандартына сәйкес ең аз пісіру уақытымен
анықталды. Макарон өнімдерінің пісіру немесе кептіруден кейінгі
механикалық қасиеттерін анықтау: беріктік, сынғыштық, оң нәтиже және
т.б. (4-сурет).

Ұннан жасалған макарон өнімдеріне арналған шикізаттың желімдеу
қасиеттерін анықтау үшін жылдам тұтқырлық анализаторы (RVA) пайдаланылды
{[}12{]}. Өнімнің құрылымы тұрақты микрожүйелік текстура анализаторы
TA-XT көмегімен анықталды. Ол макарон үлгісін кесуге қажетті күшті жазу
үшін кесу режимінде пайдаланылды. Әрбір макарон үлгісі тот баспайтын
болаттан жасалған табада бөлек дайындалған, әрбір жағдайда салмағы 25 г
үлгі алынып, 500 мл суда қайнатылған. Макарон қайнаған суға қосылып,
алдын ала белгіленген уақытқа қайнатылды. Макарон түсі, құрылымы, хош
иісі, дәмі және жалпы қолайлылығы бойынша бағаланды {[}13,14{]}.
\end{multicols}

%% \begin{longtable}[]{@{}
%%   >{\raggedright\arraybackslash}p{(\linewidth - 0\tabcolsep) * \real{1.0000}}@{}}
%% \toprule\noalign{}
%% \begin{minipage}[b]{\linewidth}\centering
%% \end{minipage} \\
%% \midrule\noalign{}
%% \endhead
%% \bottomrule\noalign{}
%% \endlastfoot
%% \end{longtable}

{\bfseries 4-сурет. Макарон өнімдерінің пісіру немесе кептіруден кейінгі
механикалық қасиеттерін анықтау: беріктік, сынғыштық, оң нәтиже және
т.б.}

{\bfseries 3-кесте. Макарон өнімдерінің су белсенділігін (aw) анықтау}

%% \begin{longtable}[]{@{}
%%   >{\centering\arraybackslash}p{(\linewidth - 8\tabcolsep) * \real{0.0874}}
%%   >{\raggedright\arraybackslash}p{(\linewidth - 8\tabcolsep) * \real{0.3185}}
%%   >{\centering\arraybackslash}p{(\linewidth - 8\tabcolsep) * \real{0.1851}}
%%   >{\centering\arraybackslash}p{(\linewidth - 8\tabcolsep) * \real{0.1360}}
%%   >{\centering\arraybackslash}p{(\linewidth - 8\tabcolsep) * \real{0.2730}}@{}}
%% \toprule\noalign{}
%% \begin{minipage}[b]{\linewidth}\centering
%% №
%% \end{minipage} & \begin{minipage}[b]{\linewidth}\centering
%% {\bfseries Үлгі атауы}
%% \end{minipage} & \begin{minipage}[b]{\linewidth}\centering
%% {\bfseries Нәтиже}
%% \end{minipage} & \begin{minipage}[b]{\linewidth}\centering
%% {\bfseries °С}
%% \end{minipage} & \begin{minipage}[b]{\linewidth}\centering
%% {\bfseries Тестілеу әдістемесі}
%% \end{minipage} \\
%% \midrule\noalign{}
%% \endhead
%% \bottomrule\noalign{}
%% \endlastfoot
%% 1 & 98,5 г ұн, 1,5г балдыркөк ұнтағы & 0,325 & 28,8 &
%% \multirow{2}{=}{\centering\arraybackslash ISO 18787:2017} \\
%% 2 & 90 г ұн, 10 г йогурт ұнтағы & 0,302 & 30,6 \\
%% \end{longtable}

\fig{p/image27}[]

{\bfseries 5 - сурет. LabSwift-aw су белсенділігі анализаторы}

3-кестеде - Макарон өнімдерінің су белсенділігін (aw) LabSwift-aw су
белсенділігі анализаторы - зертханада да, тікелей өндірісте де сұйық
және қатты үлгілердің барлық түрлерінің су белсенділігін өлшеуге
арналған портативті құрылғы көрсетілген (5-сурет). Йогурт ұнтағы
қосылған макарон өнімдерінің ылғалдылығын анықтау барысы (6-сурет)

Қолдану аймақтары: Нан және ұннан жасалған кондитерлік өнімдер; үлпек,
жарма, қытырлақ нан, чипсы; шай, кофе; ұнтақтар, дәмдеуіштер,
кептірілген өнімдер; құрғақ мал азығы; құрғақ сүт, құрғақ сүт өнімдері;
косметикалық және фармацевтикалық өнімдер.

{\bfseries Үлгі № 1 - Нәтижені талдау және интерпретация}

%% \begin{longtable}[]{@{}
%%   >{\raggedright\arraybackslash}p{(\linewidth - 4\tabcolsep) * \real{0.2549}}
%%   >{\raggedright\arraybackslash}p{(\linewidth - 4\tabcolsep) * \real{0.1369}}
%%   >{\raggedright\arraybackslash}p{(\linewidth - 4\tabcolsep) * \real{0.6082}}@{}}
%% \toprule\noalign{}
%% \begin{minipage}[b]{\linewidth}\centering
%% {\bfseries Көрсеткіш}
%% \end{minipage} & \begin{minipage}[b]{\linewidth}\centering
%% {\bfseries Мәні}
%% \end{minipage} & \begin{minipage}[b]{\linewidth}\centering
%% {\bfseries Интерпретация}
%% \end{minipage} \\
%% \midrule\noalign{}
%% \endhead
%% \bottomrule\noalign{}
%% \endlastfoot
%% Су белсенділігі (aw) & 0.325 & Өте төмен. Микробтар, ашытқылар, зең өсе
%% алмайды \\
%% Температура & 28,8°C & Өлшеу үшін қолайлы. Нақты мән алынған \\
%% Өнім ылғалдылығы & төмен & Суға бай емес, бұл сақтау тұрақтылығын
%% арттырады \\
%% \end{longtable}

{\bfseries Үлгі № 2 - Нәтижені талдау және интерпретация}

%% \begin{longtable}[]{@{}
%%   >{\raggedright\arraybackslash}p{(\linewidth - 4\tabcolsep) * \real{0.2549}}
%%   >{\centering\arraybackslash}p{(\linewidth - 4\tabcolsep) * \real{0.1369}}
%%   >{\raggedright\arraybackslash}p{(\linewidth - 4\tabcolsep) * \real{0.6082}}@{}}
%% \toprule\noalign{}
%% \begin{minipage}[b]{\linewidth}\centering
%% {\bfseries Көрсеткіш}
%% \end{minipage} & \begin{minipage}[b]{\linewidth}\centering
%% {\bfseries Мәні}
%% \end{minipage} & \begin{minipage}[b]{\linewidth}\centering
%% {\bfseries Интерпретация}
%% \end{minipage} \\
%% \midrule\noalign{}
%% \endhead
%% \bottomrule\noalign{}
%% \endlastfoot
%% Су белсенділігі (aw) & 0.302 & Өте төмен. Микробтардың ешқайсысы
%% тіршілік ете алмайды \\
%% Температура & 30,6°C & Жоғарырақ, бірақ aw өлшеуге қолайлы \\
%% Өнім ылғалдылығы & төмен & Суға бай емес, бұл сақтау тұрақтылығын
%% арттырады \\
%% \end{longtable}

{\bfseries 4-кесте - Макарон өнімдерінің ылғалдылығын анықтау}

%% \begin{longtable}[]{@{}
%%   >{\raggedright\arraybackslash}p{(\linewidth - 8\tabcolsep) * \real{0.0499}}
%%   >{\raggedright\arraybackslash}p{(\linewidth - 8\tabcolsep) * \real{0.3249}}
%%   >{\centering\arraybackslash}p{(\linewidth - 8\tabcolsep) * \real{0.2102}}
%%   >{\centering\arraybackslash}p{(\linewidth - 8\tabcolsep) * \real{0.2402}}
%%   >{\centering\arraybackslash}p{(\linewidth - 8\tabcolsep) * \real{0.1749}}@{}}
%% \toprule\noalign{}
%% \begin{minipage}[b]{\linewidth}\centering
%% №
%% \end{minipage} & \begin{minipage}[b]{\linewidth}\centering
%% Үлгі атауы
%% \end{minipage} & \begin{minipage}[b]{\linewidth}\centering
%% кептіру алдында талдауға арналған сынама бөтелкесінің салмағы (г)
%% \end{minipage} & \begin{minipage}[b]{\linewidth}\centering
%% кептіргеннен кейін талдауға арналған сынама бөтелкесінің салмағы (г)
%% \end{minipage} & \begin{minipage}[b]{\linewidth}\centering
%% талдау үшін үлгі массасы
%% 
%% (г)
%% \end{minipage} \\
%% \midrule\noalign{}
%% \endhead
%% \bottomrule\noalign{}
%% \endlastfoot
%% 1 & 98,5 г ұн, 1,5г балдыркөк ұнтағы & 1,4669 & 0,7715 & 0,7976 \\
%% 2 & 90 г ұн, 10 г йогурт ұнтағы & 1,4189 & 0,7694 & 0,7542 \\
%% \end{longtable}

\fig{p/image28}[6 - сурет. Йогурт ұнтағы қосылған макарон өнімдерінің
ылғалдылығын анықтау барысы]

\begin{multicols}{2}
{\bfseries Қорытынды.} Сонымен қорытындылай келе, Қазақстан Республикасының
аумағында макарон йогурт ұнтағының әр-түрлі пропорцияларын қолдану
арқылы дайындалды. Нәтижелерді йогурт концентрациясын жоғарылату
талшықтың құрамын арттырып, пісіру уақытын қысқартқанын, ал макарон
өнімдерінің жұмсақтығы бақылау үлгісіне қарағанда жоғарылағанын
көрсетті. Қоспаларды алу үшін шикізат ретінде йогурт, балдыркөк
ұнтақтары алынды. Бидай ұны ішінара алмастыратын 1,5\% балдыркөк ұнтағы
және 10\% ұнтағын қосу арқылы макарон өнімдерін витаминдер мен
минералдармен байыту арқылы ұлт денсаулығын жақсарту мәселелерін шешуге
көмектеседі, ресурстарды үнемдеуді (қатты бидай) қамтамасыз етеді, өнім
желісін кеңейтеді, нарықтық стандарттау мен қолжетімді бағаға байланысты
өндірістің бәсекеге қабілеттілігі мен экономикалық тиімділігін арттыруға
ықпал етеді.
\end{multicols}

\begin{center}
{\bfseries Әдебиеттер}
\end{center}

\begin{refs}
1. Толысбаев Б.С., Сейсинбинова А.А. Қазақстан Республикасындағы тамақ
өнеркәсібі кәсіпорындарының инновациялық дамуы: проблемалары мен
перспективалары// Экономика: стратегия и практика. -2019. -T.1 (14). -
Б.109-122.

2. Шаймерденова Д.А., Омаралиева А.М., Тарабаев Б.К., Сарбасова Л.Т.,
Ануарбекова С.С., Искакова Д.Б., Шаймерденов А.А., Тастанов Д.А.
Производство обогащенных макаронных изделий комплексом нанокарбоксилатов
и тонкодисперсным порошком зерновых культур// Вестник КазУТБ. -2024. -№
3(24). - С.430-442. DOI
\href{https://doi.org/10.58805/kazutb.v.3.24-407}{10.58805/kazutb.v.3.24-407}.

3. Нұрдан Д.Н. Дәстүрлі емес ұнды шикізаттан жасалған макарон өндірісі
технологиясының нормативтік қамтамасыздануы: PhD 8D07501 -"Стандарттау
және сертификаттау" мамандығы бойынша-салалар бойынша) /Алматы. - 2024.
- 172 б.

4. Искакова Г.К. Технология макаронного производства: Сырье и материалы.
Учебное пособие//Алматы: «Полиграфия-сервис и Ко». -2014. -208 с. ISBN
978-601-263-251-4

5. СТ РК ГОСТ Р 51865-2010. Макаронные изделия. Общие технические
условия/Астана. - 2010. - 56 с.

6. ГОСТ ISO 17718-2015. Зерно и мука из мягкой пшеницы. Определение
реологических свойств теста в зависимости от условий замеса и повышения
температуры/ Стандартинформ. -2015. -32 с.

7. ГОСТ 14849-89 Правила приемки и методы определения качества/ ИПК
Издательство стандартов. -1996.

8. ТР ТС 021/2011 Технический регламент Таможенного союза «О безопасности
пищевой продукции»/ЕЭК. -2011.

9. ГОСТ 31743-2017 Изделия макаронные.Общие технические
условия/Стандартинформ. -2017.

10. Bauziane M., Cappa C., Moles A., Brandolini A., Hidalg A., et al.
Effect of alkaline salts and whey protein isolate on the quality of
rice-maize gluten-free pasta//LWT. -2024. -Vol.191: 115623. DOI
10.1016/j.lwt.2023.115623.

11. ГОСТ 31964-2012 Изделия макаронные. Правила приемки и методы
определения качества/Стандартинформ. -2014.

12. Boudalia S., Gueroui Y., Boualem B., Bousbia A., Benada M., Leksir
C., Mezroua Y., Zemmouchi K.R., Saoud A., Chemmam M. Evaluation of
Physicochemical Properties and Sensory Qualities of Pasta Enriched with
Freeze-Dried Sweet Whey//~Scientia Agriculturae Bohemica. -2020. -Vol.
51(3). -P.75- 85. DOI 10.2478/sab-2020-0010.

13. Osipova G.A., Koryachkina S.Y., Koryachkin V.P., Seregina T.V.,
Zhugina A. Effects of protein-containing additives on pasta quality and
biological value//~Foods and Raw Materials. -2019. -Vol.7(1). -P.
60-66. DOI
\href{http://doi.org/10.21603/2308-4057-2019-1-60-66}{10.21603/2308-4057-2019-1-60-66}.

14. Marco M., Michela V., Carlo G.R., Erica P. Design of a Plant-Based
Yogurt-Like Product Fortified with Hemp Flour: Formulation and
Characterization// Foods. -2023. -Vol.~12 (3): 485.
\href{https://doi.org/10.3390/foods12030485}{DOI 10.3390/foods12030485}.
\end{refs}

\begin{center}
{\bfseries References}
\end{center}

\begin{refs}
1. Tolysbaev B.S., Sejsinbinova A.A. Қazaқstan Respublikasyndaғy tamaқ
өnerkәsіbі kәsіporyndarynyң innovacijalyқ damuy: problemalary men
perspektivalary// Jekonomika: strategija i praktika. -2019. -T.1 (14). -
B.109-122. {[}in Kazakh{]}

2. Shajmerdenova D.A., Omaralieva A.M., Tarabaev B.K., Sarbasova L.T.,
Anuarbekova S.S., Iskakova D.B., Shajmerdenov A.A., Tastanov D.A.
Proizvodstvo obogashhennyh makaronnyh izdelij kompleksom
nanokarboksilatov i tonkodispersnym poroshkom zernovyh
kul' tur// Vestnik KazUTB. -2024. -№ 3(24). - S.430-442.
DOI 10.58805/kazutb.v.3.24-407. {[}in Russian{]}

3. Nұrdan D.N. Dәstүrlі emes ұndy shikіzattan zhasalғan makaron өndіrіsі
tehnologijasynyң normativtіk қamtamasyzdanuy: PhD 8D07501 -"Standarttau
zhәne sertifikattau" mamandyғy bojynsha-salalar bojynsha) /Almaty. -
2024. - 172b. {[}in Kazakh{]}

4. Iskakova G.K. Tehnologija makaronnogo proizvodstva:
Syr' e i materialy. Uchebnoe posobie//Almaty:
«Poligrafija-servis i Ko». -2014. -208 s. ISBN 978-601-263-251-4. {[}in
Russian{]}

5. ST RK GOST R 51865-2010. Makaronnye izdelija. Obshhie tehnicheskie
uslovija/Astana. - 2010. - 56 s. {[}in Russian{]}

6. GOST ISO 17718-2015. Zerno i muka iz mjagkoj pshenicy. Opredelenie
reologicheskih svojstv testa v zavisimosti ot uslovij zamesa i
povyshenija temperatury/ Standartinform. -2015. -32 s. {[}in Russian{]}

7. GOST 14849-89 Pravila priemki i metody opredelenija kachestva/ IPK
Izdatel' stvo standartov. -1996. {[}in Russian{]}

8. TR TS 021/2011 Tehnicheskij reglament Tamozhennogo sojuza «O
bezopasnosti pishhevoj produkcii»/EJeK. -2011. {[}in Russian{]}

9. GOST 31743-2017 Izdelija makaronnye.Obshhie tehnicheskie
uslovija/Standartinform. -2017. {[}in Russian{]}

10. Bauziane M., Cappa C., Moles A., Brandolini A., Hidalg A., et al.
Effect of alkaline salts and whey protein isolate on the quality of
rice-maize gluten-free pasta//LWT. -2024. -Vol.191: 115623. DOI
10.1016/j.lwt.2023.115623.

11. GOST 31964-2012 Izdelija makaronnye. Pravila priemki i metody
opredelenija kachestva/Standartinform. -2014. {[}in Russian{]}

12. Boudalia S., Gueroui Y., Boualem B., Bousbia A., Benada M., Leksir
C., Mezroua Y., Zemmouchi K.R., Saoud A., Chemmam M. Evaluation of
Physicochemical Properties and Sensory Qualities of Pasta Enriched with
Freeze-Dried Sweet Whey//~Scientia Agriculturae Bohemica. -2020. -Vol.
51(3). -P.75- 85. DOI 10.2478/sab-2020-0010.

13. Osipova G.A., Koryachkina S.Y., Koryachkin V.P., Seregina T.V.,
Zhugina A. Effects of protein-containing additives on pasta quality and
biological value//~Foods and Raw Materials. -2019. -Vol.7(1). -P.
60-66. DOI
\href{http://doi.org/10.21603/2308-4057-2019-1-60-66}{10.21603/2308-4057-2019-1-60-66}.

14. Marco M., Michela V., Carlo G.R., Erica P. Design of a Plant-Based
Yogurt-Like Product Fortified with Hemp Flour: Formulation and
Characterization// Foods. -2023. -Vol.~12 (3): 485.
\href{https://doi.org/10.3390/foods12030485}{DOI 10.3390/foods12030485}.
\end{refs}

\begin{info}
\hspace{1em}\emph{{\bfseries Авторлар туралы мәлімет}}

Тоқтабек А.Ә. - PhD докторанты, ҚазҰАЗУ, Алматы, Қазақстан,
e-mail:~aijan700@mail.ru;

Мамаева Л.А. - биология ғылымдарының кандидаты, қауым.профессор,
ҚазҰАЗУ, Алматы, Қазақстан, e-mail:~laura.mamaeva@mail.ru;

Hülya Gül - доктор профессор, Сулейман Демирел университеті, Испарта,
Түркия, е-mail: hulyagul@sdu.edu.tr;

Отуншиева А.Е. - PhD докторанты, әл-Фараби атындағы ҚазҰУ,
Алматы, Қазақстан, e-mail: 03.08.1990.43@mail.ru;

Алибекова М.Д. - магистр, оқытушы., М.Әуезов атындағы Оңтүстік
Қазақстан зерттеу университеті, Шымкент, e-mail:alibekova.malika@bk.ru.

\hspace{1em}\emph{{\bfseries Information about the authors}}

Toktabek A.A. - PhD student, KazNARU, Almaty, Kazakhstan, e-mail:
aijan700@mail.ru;

Mamayeva L.A. - Candidate of Biological Sciences, Associate Professor,
KazNARU, Almaty, Kazakhstan, e-mail:~laura.mamaeva@mail.ru;

Hülya Gül - Dr. Professor, Süleyman Demirel University, Isparta,
Türkiye, е-mail: hulyagul@sdu.edu.tr;

Otunshiyeva A.E. - PhD student, Al-Farabi Kazakh National
University, Almaty, Kazakhstan, e-mail: 03.08.1990.43@mail.ru;

Alibekova M.D. - master' s student, teacher, M. Auezov
South Kazakhstan Research University, Shymkent, Kazakhstan, e-mail:
alibekova.malika@bk.ru.
\end{info}
